\documentclass[11pt]{article}
\usepackage[margin=1in]{geometry}
\usepackage{amsmath,amssymb,amsthm}
\usepackage{tikz}
\usetikzlibrary{positioning}
\newtheorem{theorem}{Theorem}
\theoremstyle{definition}
\newtheorem*{definition*}{Definition}
\newenvironment{exambox}
  {\par\medskip\noindent\begin{center}\begin{minipage}{0.92\textwidth}%
   \hrule height 1pt \vspace{4pt}\noindent\textbf{$\bigstar$ EXAM PRIORITY.} }
  {\vspace{4pt}\hrule height 1pt \end{minipage}\end{center}\medskip}

\title{CSEN 19 — Trees and Spanning Trees (June 5, 2026)}
\date{}
\author{}

\begin{document}

\maketitle

\section*{Today: Characterization of Trees and Spanning Trees}

\hrulefill

\begin{align*}
& \text{If } G \text{ is simple and connected, then } G \text{ is a tree iff } \\
& \text{for all } u,v \in V(G), \text{ there exists exactly one path between } u \text{ and } v. \\
& \text{Equivalently: } G \text{ has no cycles.}
\end{align*}

\begin{align*}
& \text{A spanning tree of a graph } G = (V,E) \text{ is a subgraph } T = (V, E') \text{ such that} \\
& T \text{ is a tree and } E' \subseteq E.
\end{align*}

\begin{exambox}
★ The professor emphasized understanding the definitions of trees and spanning trees as fundamental concepts for the exam.
\end{exambox}

[SPOKEN 0:39–1:54]
All right. Hi, everyone. This is our last class today. we're going to talk about characterization of trees and a little bit about spanning trees. So let me start by recalling. What we've defined so far. So we defined tree. To be a graph that's connected which means that it's one piece. You can go from any vertex to any other vertex using a path plus a cyclic. So this is our original definition of a tree okay. Our original definition. This is the definition that oh that's the most standard definition if you like you know. The collision. And so the question that we want to answer to that or explore to that is are there other ways of defining trees that are equivalent, that are logically equivalent. That makes sense. And there is there's actually I can think of six different ways of defining trees. Let me show you three different ways. So here's theorem.

\begin{align*}
& \text{Let } G \text{ be a graph with } n \text{ vertices and } m \text{ edges.} \\
& \text{The following are equivalent:} \\
& (1) \text{ } G \text{ is connected and acyclic} \\
& (2) \text{ } G \text{ is connected and } m=n-1 \\
& (3) \text{ } G \text{ is acyclic and } m=n-1.
\end{align*}

[SPOKEN 1:59–2:02]
We're going to let g t be a finite graph.

[SPOKEN 2:08–2:13]
With n vertices. With n vertices.

[SPOKEN 2:20–3:35]
And m edges. So this is actually the letters that you pull in graph theory use pretty often. So n usually is the number of vertices and edges. This statement is the following the conclusion by. The following statements are equivalent are logically equivalent. That's what. These are logically equivalent. So let me tell you what these three conditions are. The first one is that G is connected. basically. So this is our definition of a cheat. This is the standard definition, right? That's the first statement to G is connected. and. the number of edges is equal to n minus one.

[SPOKEN 3:35–4:46]
The number of edges is one less than the number of vertices. Three g is a cyclic no cycles a cyclic means there's no cycles on the graph. And is n minus. That's the that's the statement of the theorem. So the theorem says that this condition here condition one is logically equivalent to condition two. And condition one is also logically called the condition three. And condition two is also logically full the condition three. Does that make sense? Any of these two statements are logically equivalent. And what does that mean? It means that if g satisfies any one of them, it also satisfies the other two okay. If two satisfies any one of them, it also satisfies the other two. So if you think about what this is saying is that if you have any two you get the third one. Like if you have connected a cyclic you get minus one. If you have connected and minus one you get a cyclic.

[SPOKEN 4:46–5:58]
And if you have a cyclic and because minus one you recollect. So we want to prove this. The first one is our original definition of a tree. As I indicated here, the other two are related. So in other words, you could have used any of these three as your definition of a tree and it would be the same. You couldn’t you wouldn’t get new graphs, right. So the one of the couple of reasons why you want to show this, but one of the reasons I wanted to show you this proof, because it illustrates an idea in when you do proofs like this where you show a bunch of statements are equivalent. All you have to show. You have to show the following cycle of implications. You have to show that one implies two, and you have to shot two implies three. And then all you have to show is three plus one. That's it. You don’t have to do like all the bi conditionals. You just have to prove three statements. So strategy of proof is the following is we prove the following implications. We prove what implies two plus three. And then three implies one. That's it. That's all we have to do.

\begin{align*}
& \text{Recall:} \text{Tree = Connected + Acyclic (original definition)} \\
& \text{Theorem: Let } G \text{ be a graph with } n \text{ vertices and } m \text{ edges.} \\
& \text{The following are equivalent:} \\
& (1) \text{ } G \text{ is connected and acyclic} \\
& (2) \text{ } G \text{ is connected and } m=n-1 \\
& (3) \text{ } G \text{ is acyclic and } m=n-1.
\end{align*}

[SPOKEN 5:58–7:09]
So if you think about what's happening here we have a cycle of implications right. So one is here. This is what we're going to show. One implies two. And then I'm going to show two implies three. And then we just have to show three implies one. And that’s two means. That means everybody’s equal to everyone. Okay. Great. So basically, why is this nice? Because it means I don't have to show that to me. I don't have to make a separate proof that one implies three, because one already implies two, which implies three. I also don't have to show that three implies two minus the three implies one, which implies. So that’s a nice idea to keep in mind. All right, well let’s do it then. We just have to pursue applications. So these are three different ways of defining trees. Actually write that down. 3%. Ways. Of defining trees or characterizing trees. So this result is a type of result that mathematicians would call a characterization.

\begin{align*}
& \text{Recall:} \text{Tree = Connected + Acyclic (original definition)} \\
& \text{Theorem: Let } G \text{ be a graph with } n \text{ vertices and } m \text{ edges.} \\
& \text{The following are equivalent:} \\
& (1) \text{ } G \text{ is connected and acyclic} \\
& (2) \text{ } G \text{ is connected and } m=n-1 \\
& (3) \text{ } G \text{ is acyclic and } m=n-1.
\end{align*}

[SPOKEN 5:58–7:09]
So if you think about what's happening here we have a cycle of implications right. So one is here. This is what we're going to show. One implies two. And then I'm going to show two implies three. And then we just have to show three implies one. And that’s two means. That means everybody’s equal to everyone. Okay. Great. So basically, why is this nice? Because it means I don't have to show that to me. I don't have to make a separate proof that one implies three, because one already implies two, which implies three. I also don't have to show that three implies two minus the three implies one, which implies. So that’s a nice idea to keep in mind. All right, well let’s do it then. We just have to pursue applications. So these are three different ways of defining trees. Actually write that down. 3%. Ways. Of defining trees or characterizing trees. So this result is a type of result that mathematicians would call a characterization.
\end{document}